\section{선별 문제 풀이 I}
\label{sec:polynomial-solutions-a}

\subsection*{문제 \thechapter.4: 다항식 나눗셈}

\begin{solution}
최고차항부터 차례로 제거하면
\[
 q=X^3-X^2-2X+2,
 \qquad r=3X+3
\]
을 얻는다. 따라서
\[
 X^5-2X^4+3X^2-X+5
 =(X^3-X^2-2X+2)(X^2-X+1)+(3X+3).
\]
나머지의 차수는 $1<2$이다.
\end{solution}

\subsection*{문제 \thechapter.7: 최대공약수와 베주 표현}

\begin{solution}
\[
 f=X^4-1,\qquad g=X^3-X
\]
에 대해
\[
 f-Xg=X^2-1
\]
이고
\[
 g=X(X^2-1)
\]
이다. 따라서
\[
 \gcd(f,g)=X^2-1
\]
이며 베주 표현은
\[
 X^2-1=1\cdot f+(-X)\cdot g
\]
이다.
\end{solution}

\subsection*{문제 \thechapter.8: 몫체에서 역원}

\begin{solution}
$\F_3[X]$에서
\[
 (X+1)(X+2)=X^2+3X+2\equiv X^2+2.
\]
몫환에서는 $X^2\equiv-1\equiv2$이므로
\[
 X^2+2\equiv2+2\equiv1\pmod{X^2+1}.
\]
따라서
\[
 \overline{X+1}^{\,-1}=\overline{X+2}.
\]
\end{solution}

\subsection*{문제 \thechapter.12: $X^q-X$의 분해}

\begin{solution}
유한체 $\F_q$의 곱셈군 $\F_q^\times$의 위수는 $q-1$이다. 라그랑주 정리에 의해 모든 $a\ne0$에 대해 $a^{q-1}=1$이고, $a=0$까지 포함하면 모든 $a\in\F_q$에 대해 $a^q-a=0$이다. 따라서 $X^q-X$는 $\F_q$의 $q$개 원소를 모두 근으로 갖는다.

한편
\[
 \prod_{a\in\F_q}(X-a)
\]
도 같은 $q$개 근을 갖는 차수 $q$의 일계수다항식이다. 두 다항식의 차는 차수가 $q$보다 작으면서 $q$개의 근을 가지므로 영다항식이다. 따라서
\[
 X^q-X=\prod_{a\in\F_q}(X-a).
\]
\end{solution}

\subsection*{문제 \thechapter.14: 도함수가 영인 다항식}

\begin{solution}
$f=\sum_i a_iX^i$라 하자. $f'=0$이면 모든 $i$에 대해 $ia_i=0$이다. $p\nmid i$이면 $i\cdot1_K$는 $K$의 영이 아닌 원소이므로 역원을 갖고, 따라서 $a_i=0$이다. 그러므로 $p$의 배수인 지수만 남아
\[
 f=\sum_j a_{pj}X^{pj}=g(X^p)
\]
로 쓸 수 있다. 역으로 $f=g(X^p)$이면 각 항의 지수가 $p$의 배수이므로 형식적 미분이 모두 $0$이다.
\end{solution}

\subsection*{문제 \thechapter.16: 몫체의 기저}

\begin{solution}
$f$가 일계수이고 $\deg f=n$이므로 모든 $g\in K[X]$는 유일하게
\[
 g=qf+r,\qquad \deg r<n
\]
으로 표현된다. 따라서 몫환의 모든 원소는
\[
 a_0+a_1\alpha+\cdots+a_{n-1}\alpha^{n-1}
\]
의 꼴이다.

선형독립성을 보자. 만약
\[
 a_0+a_1\alpha+\cdots+a_{n-1}\alpha^{n-1}=0
\]
이면 $r=a_0+a_1X+\cdots+a_{n-1}X^{n-1}$가 $(f)$에 속한다. $r\ne0$이면 $f\mid r$이므로 $\deg r\ge n$이어야 하지만 $\deg r<n$이다. 따라서 $r=0$이고 모든 $a_i=0$이다.
\end{solution}

\subsection*{문제 \thechapter.18: $\Q[X]$의 인수분해를 정수계수로 내리기}

\begin{solution}
$f=gh$가 $\Q[X]$에서의 비자명한 인수분해라고 하자. 각 인수의 계수 분모의 공배수를 곱하여
\[
 g=cg_0,\qquad h=dh_0
\]
로 쓸 수 있다. 여기서 $c,d\in\Q^\times$이고 $g_0,h_0\in\Z[X]$를 원시로 택할 수 있다. 그러면
\[
 f=cd\,g_0h_0.
\]
가우스 보조정리에 의해 $g_0h_0$는 원시이다. $f$도 원시이므로 양변의 내용 비교에서 $cd=\pm1$을 얻는다. 부호를 한 인수에 흡수하면
\[
 f=(\pm g_0)h_0
\]
는 $\Z[X]$에서의 비자명한 인수분해이다.
\end{solution}
